\documentclass[12pt,a4paper]{article}
\usepackage[UTF8]{ctex}
\usepackage{amsmath,amssymb,amsthm}
\usepackage{geometry}

\geometry{margin=2.5cm}

\title{约束的独立性与矩阵秩的关系}
\author{第11章补充说明}
\date{}

\begin{document}

\maketitle

\section{问题提出}

在约束动力学中，我们经常说：
\begin{itemize}
\item "如果约束是独立的，$A(\theta)$的秩为$k$"
\item 其中$k$是约束的数量
\end{itemize}

\textbf{问题}：
\begin{enumerate}
\item 什么是独立的约束？
\item 为什么独立约束对应秩为$k$？
\item 如果约束是非独立的呢？秩会如何变化？
\end{enumerate}

\section{什么是独立的约束？}

\subsection{定义}

\textbf{线性独立的约束}：约束矩阵$A(\theta)$的$k$行是线性独立的，如果：
\begin{itemize}
\item 不存在不全为零的标量$c_1, c_2, \ldots, c_k$使得：
\begin{equation}
c_1 \cdot \text{第1行} + c_2 \cdot \text{第2行} + \cdots + c_k \cdot \text{第$k$行} = 0
\end{equation}
\item 等价地：$A(\theta)$的行向量是线性无关的
\end{itemize}

\textbf{线性非独立的约束}：如果存在不全为零的标量使得上述等式成立，则约束是线性相关的（非独立的）。

\subsection{通俗理解}

\textbf{独立约束}：
\begin{itemize}
\item 每个约束提供"新的"限制
\item 不能从其他约束推导出某个约束
\item 每个约束都是"必要的"
\end{itemize}

\textbf{非独立约束}：
\begin{itemize}
\item 某些约束可以从其他约束推导出来
\item 存在"冗余"的约束
\item 某些约束是"不必要的"
\end{itemize}

\section{矩阵秩的定义}

\subsection{秩的定义}

\textbf{矩阵的秩}：矩阵$A$的秩是：
\begin{itemize}
\item 行秩：线性独立的行的最大数量
\item 列秩：线性独立的列的最大数量
\item 行秩 = 列秩 = 秩
\end{itemize}

\textbf{对于约束矩阵$A(\theta) \in \mathbb{R}^{k \times 6}$}：
\begin{itemize}
\item 行秩：线性独立的行的数量（最多$k$）
\item 列秩：线性独立的列的数量（最多$6$）
\item 秩 = $\min(\text{行秩}, \text{列秩})$
\end{itemize}

\section{为什么独立约束对应秩为$k$？}

\subsection{定理}

\textbf{定理}：如果$A(\theta)$的$k$行是线性独立的，那么$A(\theta)$的秩为$k$。

\textbf{证明}：

由于$A(\theta) \in \mathbb{R}^{k \times 6}$有$k$行，如果这$k$行都是线性独立的，那么：
\begin{itemize}
\item 行秩 = $k$（因为有$k$个线性独立的行）
\item 列秩 $\leq 6$（因为有$6$列）
\item 秩 = $\min(\text{行秩}, \text{列秩}) = \min(k, 6)$
\end{itemize}

如果$k \leq 6$（通常情况），那么秩 = $k$。$\square$

\subsection{逆定理}

\textbf{逆定理}：如果$A(\theta)$的秩为$k$，那么$A(\theta)$的$k$行是线性独立的。

\textbf{证明}：

如果秩为$k$，那么行秩为$k$，这意味着有$k$个线性独立的行。由于$A(\theta)$只有$k$行，这$k$行都是线性独立的。$\square$

\subsection{等价表述}

\textbf{等价性}：
\begin{itemize}
\item $A(\theta)$的$k$行线性独立 $\Leftrightarrow$ $A(\theta)$的秩为$k$
\item 这是矩阵秩的基本性质
\end{itemize}

\section{如果约束是非独立的呢？}

\subsection{非独立约束的情况}

\textbf{情况}：如果$A(\theta)$的$k$行不是线性独立的，那么：
\begin{itemize}
\item 存在某些行可以从其他行线性组合得到
\item 这些行是"冗余"的
\item 实际有效的约束数量小于$k$
\end{itemize}

\textbf{秩的变化}：
\begin{itemize}
\item 如果$A(\theta)$的$k$行中只有$r$行是线性独立的（$r < k$）
\item 那么$A(\theta)$的秩为$r$（而不是$k$）
\item 实际有效的约束数量为$r$
\end{itemize}

\subsection{具体例子}

\textbf{例子1：独立约束（$k = 3$，秩 = 3）}

设$A(\theta) = \begin{bmatrix}
1 & 0 & 0 & 0 & 0 & 0 \\  % 约束1：$\omega_x = 0$
0 & 1 & 0 & 0 & 0 & 0 \\  % 约束2：$\omega_y = 0$
0 & 0 & 0 & 0 & 0 & 1     % 约束3：$v_z = 0$
\end{bmatrix}$

\textbf{检查独立性}：
\begin{itemize}
\item 行1：$\begin{bmatrix} 1 & 0 & 0 & 0 & 0 & 0 \end{bmatrix}$
\item 行2：$\begin{bmatrix} 0 & 1 & 0 & 0 & 0 & 0 \end{bmatrix}$
\item 行3：$\begin{bmatrix} 0 & 0 & 0 & 0 & 0 & 1 \end{bmatrix}$
\end{itemize}

这些行是线性独立的（没有一行可以表示为其他行的线性组合）。

\textbf{秩}：$\text{rank}(A(\theta)) = 3 = k$ ✓

\textbf{例子2：非独立约束（$k = 3$，但秩 = 2）}

设$A(\theta) = \begin{bmatrix}
1 & 0 & 0 & 0 & 0 & 0 \\  % 约束1：$\omega_x = 0$
0 & 1 & 0 & 0 & 0 & 0 \\  % 约束2：$\omega_y = 0$
1 & 1 & 0 & 0 & 0 & 0     % 约束3：$\omega_x + \omega_y = 0$（冗余！）
\end{bmatrix}$

\textbf{检查独立性}：
\begin{itemize}
\item 行1：$\begin{bmatrix} 1 & 0 & 0 & 0 & 0 & 0 \end{bmatrix}$
\item 行2：$\begin{bmatrix} 0 & 1 & 0 & 0 & 0 & 0 \end{bmatrix}$
\item 行3：$\begin{bmatrix} 1 & 1 & 0 & 0 & 0 & 0 \end{bmatrix}$ = 行1 + 行2
\end{itemize}

行3是行1和行2的线性组合，因此不是独立的。

\textbf{秩}：$\text{rank}(A(\theta)) = 2 < k = 3$

\textbf{实际有效的约束}：
\begin{itemize}
\item 只有2个独立的约束：$\omega_x = 0$和$\omega_y = 0$
\item 约束3是冗余的（可以从约束1和2推导出来）
\end{itemize}

\textbf{例子3：完全非独立约束（$k = 3$，但秩 = 1）}

设$A(\theta) = \begin{bmatrix}
1 & 0 & 0 & 0 & 0 & 0 \\  % 约束1：$\omega_x = 0$
2 & 0 & 0 & 0 & 0 & 0 \\  % 约束2：$2\omega_x = 0$（冗余！）
3 & 0 & 0 & 0 & 0 & 0     % 约束3：$3\omega_x = 0$（冗余！）
\end{bmatrix}$

\textbf{检查独立性}：
\begin{itemize}
\item 行1：$\begin{bmatrix} 1 & 0 & 0 & 0 & 0 & 0 \end{bmatrix}$
\item 行2：$\begin{bmatrix} 2 & 0 & 0 & 0 & 0 & 0 \end{bmatrix}$ = $2 \times$ 行1
\item 行3：$\begin{bmatrix} 3 & 0 & 0 & 0 & 0 & 0 \end{bmatrix}$ = $3 \times$ 行1
\end{itemize}

行2和行3都是行1的倍数，因此不是独立的。

\textbf{秩}：$\text{rank}(A(\theta)) = 1 < k = 3$

\textbf{实际有效的约束}：
\begin{itemize}
\item 只有1个独立的约束：$\omega_x = 0$
\item 约束2和3是冗余的
\end{itemize}

\section{非独立约束的影响}

\subsection{对零空间的影响}

\textbf{零空间维度}：
\begin{itemize}
\item 如果$A(\theta)$的秩为$r$（$r < k$）
\item 那么零空间的维度为$6 - r$（而不是$6 - k$）
\item 实际可以自由运动的方向更多
\end{itemize}

\textbf{例子}：
\begin{itemize}
\item 如果$k = 3$但秩$r = 2$
\item 零空间维度：$6 - 2 = 4$（而不是$6 - 3 = 3$）
\item 有4个自由度可以自由运动
\end{itemize}

\subsection{对约束力子空间的影响}

\textbf{约束力子空间维度}：
\begin{itemize}
\item 如果$A(\theta)$的秩为$r$（$r < k$）
\item 那么约束力子空间的维度为$r$（而不是$k$）
\item 实际需要控制的约束力方向更少
\end{itemize}

\textbf{例子}：
\begin{itemize}
\item 如果$k = 3$但秩$r = 2$
\item 约束力子空间维度：$2$（而不是$3$）
\item 只需要控制2个方向的约束力
\end{itemize}

\subsection{对矩阵$A\Lambda^{-1}A^T$的影响}

\textbf{可逆性}：
\begin{itemize}
\item 如果$A(\theta)$的秩为$r$（$r < k$）
\item 那么$A(\theta)\Lambda^{-1}(\theta)A^T(\theta) \in \mathbb{R}^{k \times k}$的秩为$r$
\item 这个矩阵是奇异的（不可逆）
\end{itemize}

\textbf{问题}：
\begin{itemize}
\item 在计算拉格朗日乘数时，需要计算$(A\Lambda^{-1}A^T)^{-1}$
\item 如果矩阵是奇异的，无法求逆
\item 这会导致计算失败
\end{itemize}

\textbf{解决方法}：
\begin{itemize}
\item 识别并移除冗余约束
\item 使用伪逆$(A\Lambda^{-1}A^T)^{\dagger}$代替逆
\item 或者重新定义约束矩阵，只保留独立约束
\end{itemize}

\section{如何判断约束是否独立？}

\subsection{方法1：计算秩}

\textbf{步骤}：
\begin{enumerate}
\item 计算$A(\theta)$的秩
\item 如果秩 = $k$，则约束是独立的
\item 如果秩 < $k$，则约束是非独立的，秩就是独立约束的数量
\end{enumerate}

\textbf{例子}：
\begin{itemize}
\item 如果$A(\theta) \in \mathbb{R}^{3 \times 6}$的秩为$3$，则3个约束都是独立的
\item 如果秩为$2$，则只有2个约束是独立的，1个是冗余的
\end{itemize}

\subsection{方法2：行简化}

\textbf{步骤}：
\begin{enumerate}
\item 对$A(\theta)$进行行简化（高斯消元）
\item 非零行的数量就是独立约束的数量
\item 如果非零行数量 < $k$，则存在冗余约束
\end{enumerate}

\subsection{方法3：检查线性相关性}

\textbf{步骤}：
\begin{enumerate}
\item 检查每一行是否可以表示为其他行的线性组合
\item 如果可以，则该行是冗余的
\item 否则，该行是独立的
\end{enumerate}

\section{实际应用中的处理}

\subsection{识别冗余约束}

\textbf{问题}：在实际应用中，如何识别冗余约束？

\textbf{方法}：
\begin{enumerate}
\item \textbf{计算秩}：计算$A(\theta)$的秩，如果秩 < $k$，则存在冗余约束
\item \textbf{奇异值分解}：使用SVD识别线性相关的行
\item \textbf{物理分析}：从物理角度分析哪些约束是冗余的
\end{enumerate}

\subsection{移除冗余约束}

\textbf{方法}：
\begin{enumerate}
\item 保留$r$个线性独立的行（$r$是秩）
\item 移除其他$k-r$个冗余的行
\item 重新定义约束矩阵$A_{\text{reduced}}(\theta) \in \mathbb{R}^{r \times 6}$
\end{enumerate}

\textbf{例子}：
\begin{itemize}
\item 如果$A(\theta) \in \mathbb{R}^{3 \times 6}$的秩为$2$
\item 可以移除1个冗余行，得到$A_{\text{reduced}}(\theta) \in \mathbb{R}^{2 \times 6}$
\item 使用$A_{\text{reduced}}(\theta)$进行后续计算
\end{itemize}

\subsection{使用伪逆}

\textbf{方法}：
\begin{enumerate}
\item 如果$A\Lambda^{-1}A^T$是奇异的，使用伪逆
\item 伪逆$(A\Lambda^{-1}A^T)^{\dagger}$总是存在
\item 但需要理解伪逆的物理意义
\end{enumerate}

\section{总结}

\subsection{关键结论}

\begin{enumerate}
\item \textbf{独立约束}：
\begin{itemize}
\item $A(\theta)$的$k$行线性独立 $\Leftrightarrow$ $\text{rank}(A(\theta)) = k$
\item 每个约束都是"必要的"
\end{itemize}

\item \textbf{非独立约束}：
\begin{itemize}
\item 如果某些行是线性相关的，则$\text{rank}(A(\theta)) < k$
\item 实际有效的约束数量 = $\text{rank}(A(\theta))$
\item 存在冗余约束
\end{itemize}

\item \textbf{影响}：
\begin{itemize}
\item 零空间维度：$6 - \text{rank}(A(\theta))$（而不是$6 - k$）
\item 约束力子空间维度：$\text{rank}(A(\theta))$（而不是$k$）
\item 矩阵$A\Lambda^{-1}A^T$可能奇异
\end{itemize}
\end{enumerate}

\subsection{实际建议}

\begin{itemize}
\item \textbf{设计约束时}：确保约束是独立的
\item \textbf{计算前}：检查$A(\theta)$的秩
\item \textbf{如果发现冗余}：移除冗余约束或使用伪逆
\item \textbf{物理理解}：从物理角度理解哪些约束是必要的
\end{itemize}

\subsection{关键公式}

\begin{enumerate}
\item \textbf{秩-零度定理}：
\begin{equation}
\dim(\text{Null}(A(\theta))) + \text{rank}(A(\theta)) = 6
\end{equation}

\item \textbf{如果约束独立}：
\begin{equation}
\text{rank}(A(\theta)) = k, \quad \dim(\text{Null}(A(\theta))) = 6 - k
\end{equation}

\item \textbf{如果约束非独立}：
\begin{equation}
\text{rank}(A(\theta)) = r < k, \quad \dim(\text{Null}(A(\theta))) = 6 - r > 6 - k
\end{equation}
\end{enumerate}

\end{document}

